\section{Macchine di Turing}

\begin{itemize}
    \item $M = (Q, \Sigma, \delta, s)$: Macchina di Turing deterministica.
%    \item $N = (Q, \Sigma, \nu, s)$: Macchina di Turing non deterministica.
    \item $k \in \mathbb{N}$: Numero di nastri.
\end{itemize}

\hrulefill

\begin{itemize}
    \item $Q$: Insieme finito di stati.
    \item $s \in Q$: Stato iniziale.
    \item $H = \{h, "yes", "no"\} \subset Q$: Insieme di stati terminali.
\end{itemize}

\hrulefill

\begin{itemize}
    \item $\Sigma$: Alfabeto (insieme finito di simboli) con cui viene rappresentato l'input.
    \item $\triangleright$: Start.
    \item $\sqcup$: Blank.
    \item $\Sigma_\sqcup = \Sigma \cup \{\sqcup\}$
    \item $\Gamma = \Sigma_\sqcup \cup \{\triangleright\}$ Alfabeto di lavoro che la macchina utilizza per letture e scritture.
    \item $\underline{\sigma} \in \Gamma$: Simbolo sotto la testina di un nastro.
\end{itemize}

\hrulefill

\begin{itemize}
    \item $\delta: Q \times \Gamma^k \rightarrow Q \times \Gamma^k \times D^k$ funzione di transizione: prende in input i simboli sotto le testine e lo stato corrente, restituisce simboli che sovrascrivono quelli letti, uno stato in cui transitare al passo successivo e gli spostamenti che ogni testina deve effettuare effettuare alla fine.
    \item $D = \{\leftarrow, \rightarrow, -\}$ spostamenti della testina sul nastro: puó muoversi di una cella a sinistra o a destra, oppure restare ferma.
    \item $\delta(q, [\dots, \triangleright_i, \dots]) = (q', [\dots, \triangleright_i, \dots], [\dots, \rightarrow_i, \dots])$ (con $q, q' \in Q$): se una qualsiasi testina $i$ si trova su $\triangleright$, deve riscriverlo e spostarsi obbligatoriamente a destra.
\end{itemize}

\hrulefill

\begin{itemize}
    \item $x \in \Sigma^*$: Input scritto sul primo nastro all'inizio della computazione.
    \item $n = |x| \in \mathbb{N}$
    \item $f(n)$: Tempo utilizzato dalla computazione.
    \item $g(n)$: Spazio di memoria utilizzato dalla computazione.
\end{itemize}

% \begin{itemize}
%     \item $U$: Macchina di Turing Universale
%     \item $\underline{M}$: Codifica di $M$
%     \item $X = (\underline{M}, x)$: Input di $U$
%     \item $N = |X|$
% \end{itemize}


\subsection{Configurazione}
Siano, $\forall i \in \{1, \dots, k\}$:
\begin{itemize}
    \item $q \in Q$: Stato corrente.
    \item $w_i = "\triangleright w_{i,2} w_{i,3} \dots \underline{w_{i,{m_i}}}" \in \{\triangleright\} \times {\Sigma_\sqcup}^*$: Stringa sul nastro $i$ prima della testina, incluso il simbolo sotto la testina.
    (Nota: all'inizio della computazione, $w_i = \underline{\triangleright}$).
    \item $u_i = "u_{i,1} u_{i,2} \dots u_{i,{p_i}}" \in {\Sigma_\sqcup}^*$: Stringa sul nastro $i$ dopo la testina.
\end{itemize}

\[
\configurazione{q}{w}{u}
\]

\subsubsection{Massimo numero di configurazioni possibili}
\begin{itemize}
    \item Numero di stati: $|Q|$
    \item Massimo numero di stringhe possibili: $|\Sigma_\sqcup|^{g(n) \cdot k}$ \\
    ($|\Sigma_\sqcup|^{g(n)}$ stringhe per ogni nastro)
    \item Massimo numero di posizioni delle testine possibili: $g(n)^k$ \\
    ($g(n)$ posizioni per ogni testina)
\end{itemize}

{\[|Q| \cdot |{\Sigma_\sqcup|}^{g(n) \cdot k} \cdot g(n)^k\]}


\subsection{Computazione}
\subsubsection{Computazione di un singolo passo}
\[
\configurazione{q}{w}{u}
\xrightarrow{M}
\configurazione{q'}{w'}{u'}
\]

Se e solo se: \\

Se $\delta(q, [\underline{w}]) = (q', [\sigma], [d])$, dove:
\begin{itemize}
	\item $[\underline{w}] = (\underline{w_{1,m_1}}, \dots, \underline{w_{k,m_k}})$: Simboli sotto le testine dei $k$ nastri.
    \item $q' \in Q$: Stato successivo della macchina.
    \item $[\sigma] = [\sigma_1, \dots, \sigma_k] \in {\Sigma_\sqcup}^k$: Vettore di $k$ caratteri, ognuno da scrivere su una delle $k$ celle sotto una testina.
    \item $[d] = [d_1, \dots, d_k] \in D^k$: Vettore di $k$ direzioni che ogni testina deve seguire.
\end{itemize}

Allora, $\forall i \in \{1, \dots, k\} :$
{\small
\begin{enumerate}
    \item $(d_i = -) \implies ((w'_i = "\triangleright w_{i,2} w_{i,3} \dots \underline{\sigma_i}") \land (u'_i = u_i))$
    \item $(d_i = \leftarrow) \implies ((w'_i = "\triangleright w_{i,2} w_{i,3} \dots \underline{w_{i,m_{i-1}}}") \land (u'_i = "\sigma_i u_{i,1} u_{i,2} \dots u_{i,p_i}"))$
    \item $((d_i = \rightarrow) \land (u_i \neq \epsilon)) \implies ((w'_i = "\triangleright w_{i,2} w_{i,3} \dots \sigma_i \underline{u_{i,1}}") \land (u'_i = "u_{i,2} \dots u_{i,p_i}")) $
    \item $((d_i = \rightarrow) \land (u_i = \epsilon)) \implies ((w'_i = "\triangleright w_{i,2} w_{i,3} \dots \sigma_i \underline{\sqcup}") \land (u'_i = \epsilon)) $
\end{enumerate}
}

\subsubsection{Computazione di $t$ passi}
\[
\configurazione{q_0}{w_0}{u_0}
\xrightarrow{M^t}
\configurazione{q_t}{w_t}{u_t}
\] \\

Se $\forall j \in \{0, \dots, t-1\} :$
\[
\configurazione{q_j}{w_j}{u_j}
\xrightarrow{M}
\configurazione{q_{j+1}}{w_{j+1}}{u_{j+1}}
\]

\subsubsection{Computazione di un numero di passi indefinito}
\[
\configurazione{q}{w}{u}
\xrightarrow{M^*}
\configurazione{q'}{w'}{u'}
\] \\

Se $\exists t \in \mathbb{N} :$
\[
\configurazione{q}{w}{u}
\xrightarrow{M^t}
\configurazione{q'}{w'}{u'}
\]

\subsection{Terminazione}
    $M(x) \in \Sigma^* \cup \{"yes", "no",\nearrow$\}
    \begin{itemize}
        \item $M(x) = "yes"$ sse $\exists w, u :$
        \[
        \configurazionestart{x} \xrightarrow{M^*} \configurazione{"yes"}{w}{u}
        \]
        \item $M(x) = "no"$ sse $\exists w, u :$
        \[
        \configurazionestart{x} \xrightarrow{M^*} \configurazione{"no"}{w}{u}
        \]
        \item $M(x) = y \in \Sigma^*$ sse $\exists w, u :$
        \[
        \configurazionestart{x} \xrightarrow{M^*} \configurazione{h}{w}{u} \qquad \land \qquad 
        y = "w_k u_k"
        \]
        \item $M(x) = \nearrow$ se la macchina diverge (la computazione non termina). Cioé se non vale nessuno dei casi precedenti.
	\end{itemize}

\subsection{Macchine di Turing con input e output}
\begin{itemize}
	\item $k >= 2$
	\item Il primo nastro non viene mai modificato (read only).
	\item L'ultimo nastro viene scandito solo verso destra (write only).
	\item Il cursore del primo nastro non supera mai il primo $\sqcup$ successivo all'input $x$.
\end{itemize}

\subsection{Composizione di Macchine di Turing}
Date due MdT $M_1 = (Q_1, \Sigma, \delta_1, s_1)$ e $M_2 = (Q_2, \Sigma, \delta_2, s_2)$ tali che:
\begin{itemize}
	\item Hanno lo stesso alfabeto $\Sigma$ 
	\subitem (Se gli alfabeti sono diversi, basta considerare la loro unione).
	\item $Q_1 \cap Q_2 = \emptyset$ 
	\subitem (Se gli insiemi di stati non sono disgiunti, basta rinominare gli stati di uno dei due).
	\item $M_1$ termina con tutte le testine su $\triangleright$ 
	\subitem (Se non accade, basta cambiare la funzione di transizione $\delta_1$, redirezionando le coppie che portano verso gli stati $"yes"$ e $h$ in uno stato che riporta le testine su $\triangleright$ e poi termina).
\end{itemize}

La loro \textbf{composizione} é: $(M_1, M_2) =  (Q, \Sigma, \delta, s)$, dove: 
\begin{itemize}
	\item $Q = Q_1 \cup Q_2$
	\item $s = s_1$
	\item $\forall \sigma \in \Sigma^k$:
	\[
	\delta(q, \sigma) =
	\begin{cases} 
		\delta_2(q, \sigma)  \quad se \quad q \in Q_2 \\
		\delta_1(q, \sigma) \quad se \quad q \in Q_1 \land \delta_1(q, \sigma)  = (p, \tau, d) \land p \notin \{h, "yes"\} \\
		(s_2, \tau, d) \quad se \quad q \in Q_1 \land  \delta_1(q, \sigma)  = (p, \tau, d) \land p \in \{h, "yes"\} \\
	\end{cases}
	\]
	Dove $\tau \in \Sigma^k$ \quad $d \in \{\leftarrow, \rightarrow, -\}^k$
\end{itemize}

\begin{itemize}
	\item La funzione di transizione modifica solo $\delta_1$, reindirizzando le coppie $(q, \sigma)$ che portano a uno stato terminale $h$ o $"yes"$ verso lo stato iniziale di $M_2$. 
	\item Tutti gli altri casi restano uguali alle funzioni di transizione delle due macchine.
	\item Le coppie $(q, \sigma)$ che per $\delta_1$ vanno nello stato $"no"$ non devono essere reindirizzate: \textbf{se $\mathbf{M_1}$ rigetta l'input, la computazione non deve continuare}.
\end{itemize}
